\documentclass[12pt,a4paper]{article}

% ---- Encoding & Font ----
\usepackage{iftex}
\ifPDFTeX
  \usepackage[utf8]{inputenc}
  \usepackage[T1]{fontenc}
\else
  \usepackage{fontspec}  % XeTeX/LuaTeX (tectonic): Unicode nativo (·, —, ¿, ª, §)
\fi
\usepackage[spanish]{babel}
\usepackage{csquotes}

% ---- Page layout ----
\usepackage[top=2.5cm, bottom=2.5cm, left=2.5cm, right=2.5cm]{geometry}
\usepackage{setspace}
\onehalfspacing

% ---- Math ----
\usepackage{amsmath, amssymb, amsthm}
\usepackage{mathtools}
\usepackage{physics}

% ---- Graphics ----
\usepackage{graphicx}
\usepackage{tikz}
\usepackage{float}
\usepackage{caption}

% ---- Tables ----
\usepackage{array}
\usepackage{booktabs}
\usepackage{tabularx}
\usepackage{colortbl}

% ---- Colours ----
\usepackage{xcolor}
\definecolor{notebg}{HTML}{FFF3CD}
\definecolor{noteborder}{HTML}{FFC107}
\definecolor{boxred}{HTML}{F8D7DA}
\definecolor{boxredborder}{HTML}{F5C6CB}

% ---- Boxes ----
\usepackage{mdframed}
\usepackage{tcolorbox}
\tcbuselibrary{most}

\newtcolorbox{keybox}[1][]{
  colback=blue!5,
  colframe=blue!70!black,
  arc=4pt,
  boxrule=1pt,
  fonttitle=\bfseries,
  title={#1}
}

\newtcolorbox{warnbox}[1][]{
  colback=boxred,
  colframe=boxredborder,
  coltitle=black!75,
  arc=4pt,
  boxrule=1pt,
  fonttitle=\bfseries,
  title={#1}
}

\newtcolorbox{notebox}[1][]{
  colback=notebg,
  colframe=noteborder,
  coltitle=black!75,
  arc=4pt,
  boxrule=1pt,
  fonttitle=\bfseries,
  title={#1}
}


% ---- Pseudocódigo ----
\usepackage[spanish,onelanguage,ruled,vlined]{algorithm2e}
\SetKwInput{KwEntrada}{Entrada}
\SetKwInput{KwSalida}{Salida}
% ---- Diagramas (gráficas) ----
\usepackage{pgfplots}
\pgfplotsset{compat=1.18}
\usetikzlibrary{babel}
% courses/MetNum2023/Clases/assets/tikzstyles.tex
% ---------------------------------------------------------------------------
% Estilos compartidos para los diagramas del curso Métodos Numéricos
% (MetNum2023). Fuente única del "look" visual (colores, grosores, escalas)
% para que todas las figuras (matrices triangulares, curvas de convergencia
% / error, splines a trozos) sean consistentes entre clases.
%
% Uso: la clase debe cargar antes `tikz` y `pgfplots` (bloque §6.3 del
%      CLAUDE.md del curso), y luego
%      % courses/MetNum2023/Clases/assets/tikzstyles.tex
% ---------------------------------------------------------------------------
% Estilos compartidos para los diagramas del curso Métodos Numéricos
% (MetNum2023). Fuente única del "look" visual (colores, grosores, escalas)
% para que todas las figuras (matrices triangulares, curvas de convergencia
% / error, splines a trozos) sean consistentes entre clases.
%
% Uso: la clase debe cargar antes `tikz` y `pgfplots` (bloque §6.3 del
%      CLAUDE.md del curso), y luego
%      % courses/MetNum2023/Clases/assets/tikzstyles.tex
% ---------------------------------------------------------------------------
% Estilos compartidos para los diagramas del curso Métodos Numéricos
% (MetNum2023). Fuente única del "look" visual (colores, grosores, escalas)
% para que todas las figuras (matrices triangulares, curvas de convergencia
% / error, splines a trozos) sean consistentes entre clases.
%
% Uso: la clase debe cargar antes `tikz` y `pgfplots` (bloque §6.3 del
%      CLAUDE.md del curso), y luego
%      \input{../assets/tikzstyles.tex}
% (compilar desde el directorio de la clase para que la ruta relativa resuelva).
%
% Paleta propia de este curso (no se reusa la de Fisica3-2015/CDIV2017/
% ElecMag2024): mismos tonos base para alinear con keybox/notebox/warnbox,
% pero archivo independiente, sin dependencia cruzada entre cursos.
% ---------------------------------------------------------------------------

% ---- Paleta (alineada con las cajas del preámbulo: keybox/notebox/warnbox) ----
\definecolor{figblue}{HTML}{1F4E79}   % azul principal (L, curvas principales, keybox)
\definecolor{figred}{HTML}{B03A2E}    % rojo de énfasis (U, warnbox, error)
\definecolor{figamber}{HTML}{B8860B}  % ámbar (notebox, pivotes)
\definecolor{figgray}{HTML}{555555}   % auxiliares, ejes, ceros de matriz

% ---- TikZ: librerías y estilos reutilizables ----
% Nota: las puntas se declaran como `-{Latex[...]}` y no como `->`. La forma
% con llaves NO contiene un `>` literal, así que es inmune al `>` activo que
% introduce babel español (de todos modos se carga \usetikzlibrary{babel} en
% el bloque §6.3 para cubrir cualquier `->` suelto).
\usetikzlibrary{arrows.meta,calc,matrix,patterns}

\tikzset{
  figline/.style={figblue, line width=0.9pt},
  figaux/.style={figgray, dashed, line width=0.6pt},
  figaxis/.style={figgray, line width=0.7pt, -{Latex[length=2.2mm,width=1.5mm]}},
  % estructura de matrices (L triangular inferior, U triangular superior)
  matcelda/.style={minimum width=7mm, minimum height=7mm, anchor=center, font=\small},
  matnz/.style={matcelda, fill=figblue!12},
  matnzU/.style={matcelda, fill=figred!12},
  matcero/.style={matcelda, text=figgray!70},
  matpivote/.style={matcelda, fill=figamber!25, draw=figamber, line width=0.8pt},
  % nodos de interpolación / splines
  fignodo/.style={circle, inner sep=0pt, minimum size=4.5pt, draw=figblue, line width=0.9pt, fill=white},
  fignodoactivo/.style={fignodo, fill=figamber},
  % etiquetas
  figlbl/.style={font=\small},
  figlblaux/.style={font=\footnotesize, figgray},
}

% ---- pgfplots: estilos base ----
\pgfplotsset{
  % convergencia / error vs. n (grado, iteración) — suele ir en escala log
  errorfig/.style={
    width=10cm, height=6.5cm,
    xlabel style={font=\small},
    ylabel style={font=\small},
    tick label style={font=\footnotesize},
    every axis plot/.append style={line width=1.1pt, mark=*, mark size=1.5pt},
    grid=major, grid style={figgray!15},
    legend cell align=left,
    legend style={font=\footnotesize, draw=figgray!40, at={(0.5,-0.25)}, anchor=north},
    clip=false,
  },
  % interpolación / splines: un color por tramo, nodos marcados aparte
  interpfig/.style={
    width=10cm, height=6cm,
    axis lines=middle,
    xlabel style={font=\small, at={(axis description cs:1,0.5)}, anchor=west},
    ylabel style={font=\small, at={(axis description cs:0.5,1)}, anchor=south},
    tick label style={font=\footnotesize},
    every axis plot/.append style={line width=1.1pt},
    grid=none,
    legend cell align=left,
    legend style={font=\footnotesize, draw=figgray!40},
    clip=false,
  },
}

% (compilar desde el directorio de la clase para que la ruta relativa resuelva).
%
% Paleta propia de este curso (no se reusa la de Fisica3-2015/CDIV2017/
% ElecMag2024): mismos tonos base para alinear con keybox/notebox/warnbox,
% pero archivo independiente, sin dependencia cruzada entre cursos.
% ---------------------------------------------------------------------------

% ---- Paleta (alineada con las cajas del preámbulo: keybox/notebox/warnbox) ----
\definecolor{figblue}{HTML}{1F4E79}   % azul principal (L, curvas principales, keybox)
\definecolor{figred}{HTML}{B03A2E}    % rojo de énfasis (U, warnbox, error)
\definecolor{figamber}{HTML}{B8860B}  % ámbar (notebox, pivotes)
\definecolor{figgray}{HTML}{555555}   % auxiliares, ejes, ceros de matriz

% ---- TikZ: librerías y estilos reutilizables ----
% Nota: las puntas se declaran como `-{Latex[...]}` y no como `->`. La forma
% con llaves NO contiene un `>` literal, así que es inmune al `>` activo que
% introduce babel español (de todos modos se carga \usetikzlibrary{babel} en
% el bloque §6.3 para cubrir cualquier `->` suelto).
\usetikzlibrary{arrows.meta,calc,matrix,patterns}

\tikzset{
  figline/.style={figblue, line width=0.9pt},
  figaux/.style={figgray, dashed, line width=0.6pt},
  figaxis/.style={figgray, line width=0.7pt, -{Latex[length=2.2mm,width=1.5mm]}},
  % estructura de matrices (L triangular inferior, U triangular superior)
  matcelda/.style={minimum width=7mm, minimum height=7mm, anchor=center, font=\small},
  matnz/.style={matcelda, fill=figblue!12},
  matnzU/.style={matcelda, fill=figred!12},
  matcero/.style={matcelda, text=figgray!70},
  matpivote/.style={matcelda, fill=figamber!25, draw=figamber, line width=0.8pt},
  % nodos de interpolación / splines
  fignodo/.style={circle, inner sep=0pt, minimum size=4.5pt, draw=figblue, line width=0.9pt, fill=white},
  fignodoactivo/.style={fignodo, fill=figamber},
  % etiquetas
  figlbl/.style={font=\small},
  figlblaux/.style={font=\footnotesize, figgray},
}

% ---- pgfplots: estilos base ----
\pgfplotsset{
  % convergencia / error vs. n (grado, iteración) — suele ir en escala log
  errorfig/.style={
    width=10cm, height=6.5cm,
    xlabel style={font=\small},
    ylabel style={font=\small},
    tick label style={font=\footnotesize},
    every axis plot/.append style={line width=1.1pt, mark=*, mark size=1.5pt},
    grid=major, grid style={figgray!15},
    legend cell align=left,
    legend style={font=\footnotesize, draw=figgray!40, at={(0.5,-0.25)}, anchor=north},
    clip=false,
  },
  % interpolación / splines: un color por tramo, nodos marcados aparte
  interpfig/.style={
    width=10cm, height=6cm,
    axis lines=middle,
    xlabel style={font=\small, at={(axis description cs:1,0.5)}, anchor=west},
    ylabel style={font=\small, at={(axis description cs:0.5,1)}, anchor=south},
    tick label style={font=\footnotesize},
    every axis plot/.append style={line width=1.1pt},
    grid=none,
    legend cell align=left,
    legend style={font=\footnotesize, draw=figgray!40},
    clip=false,
  },
}

% (compilar desde el directorio de la clase para que la ruta relativa resuelva).
%
% Paleta propia de este curso (no se reusa la de Fisica3-2015/CDIV2017/
% ElecMag2024): mismos tonos base para alinear con keybox/notebox/warnbox,
% pero archivo independiente, sin dependencia cruzada entre cursos.
% ---------------------------------------------------------------------------

% ---- Paleta (alineada con las cajas del preámbulo: keybox/notebox/warnbox) ----
\definecolor{figblue}{HTML}{1F4E79}   % azul principal (L, curvas principales, keybox)
\definecolor{figred}{HTML}{B03A2E}    % rojo de énfasis (U, warnbox, error)
\definecolor{figamber}{HTML}{B8860B}  % ámbar (notebox, pivotes)
\definecolor{figgray}{HTML}{555555}   % auxiliares, ejes, ceros de matriz

% ---- TikZ: librerías y estilos reutilizables ----
% Nota: las puntas se declaran como `-{Latex[...]}` y no como `->`. La forma
% con llaves NO contiene un `>` literal, así que es inmune al `>` activo que
% introduce babel español (de todos modos se carga \usetikzlibrary{babel} en
% el bloque §6.3 para cubrir cualquier `->` suelto).
\usetikzlibrary{arrows.meta,calc,matrix,patterns}

\tikzset{
  figline/.style={figblue, line width=0.9pt},
  figaux/.style={figgray, dashed, line width=0.6pt},
  figaxis/.style={figgray, line width=0.7pt, -{Latex[length=2.2mm,width=1.5mm]}},
  % estructura de matrices (L triangular inferior, U triangular superior)
  matcelda/.style={minimum width=7mm, minimum height=7mm, anchor=center, font=\small},
  matnz/.style={matcelda, fill=figblue!12},
  matnzU/.style={matcelda, fill=figred!12},
  matcero/.style={matcelda, text=figgray!70},
  matpivote/.style={matcelda, fill=figamber!25, draw=figamber, line width=0.8pt},
  % nodos de interpolación / splines
  fignodo/.style={circle, inner sep=0pt, minimum size=4.5pt, draw=figblue, line width=0.9pt, fill=white},
  fignodoactivo/.style={fignodo, fill=figamber},
  % etiquetas
  figlbl/.style={font=\small},
  figlblaux/.style={font=\footnotesize, figgray},
}

% ---- pgfplots: estilos base ----
\pgfplotsset{
  % convergencia / error vs. n (grado, iteración) — suele ir en escala log
  errorfig/.style={
    width=10cm, height=6.5cm,
    xlabel style={font=\small},
    ylabel style={font=\small},
    tick label style={font=\footnotesize},
    every axis plot/.append style={line width=1.1pt, mark=*, mark size=1.5pt},
    grid=major, grid style={figgray!15},
    legend cell align=left,
    legend style={font=\footnotesize, draw=figgray!40, at={(0.5,-0.25)}, anchor=north},
    clip=false,
  },
  % interpolación / splines: un color por tramo, nodos marcados aparte
  interpfig/.style={
    width=10cm, height=6cm,
    axis lines=middle,
    xlabel style={font=\small, at={(axis description cs:1,0.5)}, anchor=west},
    ylabel style={font=\small, at={(axis description cs:0.5,1)}, anchor=south},
    tick label style={font=\footnotesize},
    every axis plot/.append style={line width=1.1pt},
    grid=none,
    legend cell align=left,
    legend style={font=\footnotesize, draw=figgray!40},
    clip=false,
  },
}


% ---- Hyperlinks & TOC ----
\usepackage[hidelinks]{hyperref}
\usepackage{bookmark}

% ---- Enumerate ----
\usepackage{enumitem}

% ---- Miscellaneous ----
\usepackage{icomma}
\usepackage{siunitx}
\usepackage{textcomp}
\usepackage[bottom]{footmisc}

% ---- Title ----
\title{\textbf{Clase 11: Métodos Numéricos} \\[0.3em]
        \large{Criterios de parada e interpolación polinomial}}
\author{Transcripto y expandido --- Curso Métodos Numéricos (OpenFING)}
\date{}

\begin{document}
\maketitle
\thispagestyle{empty}
\newpage
\tableofcontents
\newpage

\section{Cuándo detener una iteración}

\subsection{Parar por residuo normalizado}

La clase cierra sistemas lineales preguntando cuándo detener un método iterativo. Se mencionan Jacobi, Gauss–Seidel y variantes relajadas, ya tratadas en clases anteriores. La discusión usa el esquema

\[
\mathbf x^{(k+1)}=Q\mathbf x^{(k)}+\mathbf c,
\qquad \mathbf x^*=Q\mathbf x^*+\mathbf c.
\]

Aquí $\mathbf c$ designa el término fijo del método para no confundirlo con el residuo $\mathbf r^{(k)}=A\mathbf x^{(k)}-\mathbf b$. La transcripción usa la misma letra para ambos y los distingue con colores.

\begin{notebox}[Alcance y precisiones]
Las clases 9 y 10 no forman parte de este lote. Se utilizan las relaciones de iteración y error recordadas en esta clase, sin reconstruir el desarrollo de los métodos anteriores.
\end{notebox}

Una primera propuesta es detenerse cuando la norma del residuo es menor que una tolerancia prefijada. Para dar una interpretación relativa, el docente modifica el criterio a

\[
\boxed{\|\mathbf r^{(k)}\|<\varepsilon\|\mathbf b\|}.
\]

Se supone $\mathbf b\ne0$ y $A$ invertible. La norma matricial empleada en $\kappa(A)$ es la inducida por la norma vectorial elegida. Recordando la estimación de error y residuo:

\[
\frac{\|\mathbf x^{(k)}-\mathbf x^*\|}{\|\mathbf x^*\|}
\le\kappa(A)\frac{\|\mathbf r^{(k)}\|}{\|\mathbf b\|}
<\kappa(A)\varepsilon.
\]

El residuo es computable y la norma de $\mathbf b$ ya se conoce. La tolerancia no controla directamente el error relativo: interviene el número de condición. Qué tan grande es tolerable ese factor depende de la precisión que se necesite. La clase no fija una tolerancia universal.

\subsection{Parar porque los iterados cambian poco}

Otra propuesta es detenerse cuando

\[
\|\mathbf x^{(k+1)}-\mathbf x^{(k)}\|<\varepsilon.
\]

El hecho de moverse poco parece indicar proximidad a la solución, pero hay que justificarlo. Sea $\mathbf e^{(k)}=\mathbf x^{(k)}-\mathbf x^*$. Restando las relaciones del método y del punto fijo se obtiene $\mathbf e^{(k+1)}=Q\mathbf e^{(k)}$. Se suma y resta el siguiente iterado para hacer aparecer exactamente la cantidad controlada:

\[
\begin{aligned}
\mathbf e^{(k)}
&=\mathbf x^{(k)}-\mathbf x^{(k+1)}
  +\mathbf x^{(k+1)}-\mathbf x^*\\
&=\mathbf x^{(k)}-\mathbf x^{(k+1)}+Q\mathbf e^{(k)}.
\end{aligned}
\]

La desigualdad triangular debe agrupar la diferencia de iterados, no separar sus normas individuales. Después se aplica compatibilidad:

\[
\|\mathbf e^{(k)}\|
\le\|\mathbf x^{(k)}-\mathbf x^{(k+1)}\|
+\|Q\|\|\mathbf e^{(k)}\|.
\]

Si $\|Q\|<1$, se puede pasar el último término a la izquierda y dividir por un número positivo:

\[
\boxed{\|\mathbf e^{(k)}\|
\le\frac{\|\mathbf x^{(k+1)}-\mathbf x^{(k)}\|}{1-\|Q\|}
<\frac{\varepsilon}{1-\|Q\|}}.
\]

La cota corresponde al error del iterado $k$ y es absoluta. Si $\|Q\|$ está cerca de uno, el factor es grande; si $\|Q\|\ge1$, este despeje no da el control buscado. Puede haber un método convergente para el que la norma elegida no satisfaga $\|Q\|<1$. El docente recuerda la relación de la convergencia con el radio espectral, pero no demuestra aquí la existencia de una norma adecuada ni un resultado sobre ínfimos de normas.

Para comparar con un criterio relativo se podría escalar la tolerancia por una norma del iterado, observación mencionada sin desarrollar. No se añade una garantía relativa a la desigualdad absoluta obtenida.

\subsection{El límite de iteraciones no es convergencia}

Además del criterio de precisión, se fija una cantidad máxima de iteraciones para evitar que el programa continúe indefinidamente. Alcanzar ese límite sólo significa detener la ejecución, no haber alcanzado la tolerancia. Una transcripción algorítmica del criterio de residuo es:

\begin{algorithm}[H]
\small\DontPrintSemicolon
\hspace*{0.0em}\texttt{Entrada: A, b, regla de iteracion, x inicial, epsilon, Kmax}\;
\hspace*{0.0em}\texttt{Para k = 0, ..., Kmax:}\;
\hspace*{0.7em}\texttt{r = A*x - b}\;
\hspace*{0.7em}\texttt{Si norma(r) < epsilon * norma(b):}\;
\hspace*{1.4em}\texttt{devolver x con tolerancia alcanzada}\;
\hspace*{0.7em}\texttt{Si k = Kmax:}\;
\hspace*{1.4em}\texttt{devolver x con limite alcanzado}\;
\hspace*{0.7em}\texttt{x = siguiente\_iterado(x)}\;
\end{algorithm}

Se calcula el residuo del valor que se va a devolver. Sólo se actualiza si todavía no se alcanzó la tolerancia y queda presupuesto de iteraciones. Los nombres y la distinción de salidas explicitan editorialmente la red de seguridad propuesta; no son un código completo del método iterativo.

En aritmética exacta, si $\mathbf x^{(k)}\to\mathbf x^*$, el residuo tiende a cero. Esa afirmación no proporciona una cantidad práctica de iteraciones ni elimina las limitaciones de precisión finita. El docente menciona las matrices de Hilbert del práctico como ejemplos problemáticos, sin desarrollar aquí sus entradas ni realizar una prueba de tiempo.

\section{Planteo de la interpolación}

\subsection{Unir puntos no determina una función}

Dados $n+1$ puntos $(x_i,y_i)$, con $i=0,\ldots,n$ y abscisas distintas dos a dos, se busca una función real $\phi$ tal que $\phi(x_i)=y_i$. Las abscisas deben ser distintas porque dos alturas diferentes sobre una misma abscisa no pueden pertenecer al gráfico de una función. Las ordenadas sí pueden coincidir.

Sin elegir una clase de funciones hay infinitas soluciones. La clase restringe el problema a \textbf{polinomios algebraicos}. Se menciona la interpolación trigonométrica, útil para señales periódicas, como otro problema con herramientas diferentes, sin desarrollarla.

Los datos discretos pueden usarse para construir una función que luego se integre, derive o cuyas raíces se busquen. La interpolación conecta información puntual con problemas continuos. También se mencionan diseño asistido por computadora, curvas y letras vectoriales: representar una forma mediante puntos y reglas geométricas permite escalarla sin ampliar píxeles.

\begin{notebox}[Alcance y precisiones]
El dibujo de una letra y sus puntos de control se presenta como motivación. El audio no determina su contorno ni las coordenadas; no se reconstruye una tipografía concreta ni se afirma que todos los puntos de control sean necesariamente puntos de interpolación.
\end{notebox}

\subsection{El grado correcto es una cota}

Un punto determina una constante. Dos puntos con abscisas distintas determinan una recta, que puede ser horizontal. Por eso la condición es grado \textbf{menor o igual} a uno, no grado exactamente uno. En general:

\[
\boxed{\begin{gathered}
x_i\ne x_j\ (i\ne j)\quad\Longrightarrow\quad
\exists!\ p_n\in\mathcal P_n,\\
p_n(x_i)=y_i\quad(i=0,\ldots,n),
\end{gathered}}
\]

 donde $\mathcal P_n$ es el espacio de polinomios de grado a lo sumo $n$. Hay $n+1$ datos y $n$ es la cota de grado. El nombre \textbf{polinomio interpolante} se usa con esa restricción, que hace único el problema. Si se permiten grados mayores, se pierde esa unicidad.

\section{Existencia y unicidad por construcción}

\subsection{Caso base y paso inductivo}

La prueba hace existencia y unicidad conjuntamente por inducción. Para $n=0$ sólo hay un punto $(x_0,y_0)$ y la única constante apropiada es $p_0(x)=y_0$.

Supóngase el resultado para $n$ puntos y grado a lo sumo $n-1$. Entre los $n+1$ puntos actuales, se toman los primeros $n$, de índice 0 a $n-1$. Por hipótesis hay un único $p_{n-1}$ que los interpola. Falta incorporar $(x_n,y_n)$.

Se busca $p_n=p_{n-1}+q$, con $\deg q\le n$. Esto no restringe artificialmente la búsqueda: cualquier candidato $p_n\in\mathcal P_n$ puede escribirse así, definiendo $q=p_n-p_{n-1}$. Para no alterar los puntos ya correctos, debe cumplirse

\[
q(x_i)=0,\qquad i=0,\ldots,n-1.
\]

Un polinomio de grado a lo sumo $n$ con esas $n$ raíces distintas sólo puede tener la forma

\[
q(x)=a_n\prod_{i=0}^{n-1}(x-x_i).
\]

La libertad quedó reducida a una constante $a_n$. Se permite $a_n=0$: si el nuevo dato ya está sobre $p_{n-1}$, no se necesita subir efectivamente el grado.

\subsection{Ajustar el nuevo punto}

Se exige $p_n(x_n)=y_n$:

\[
y_n=p_{n-1}(x_n)+a_n\prod_{i=0}^{n-1}(x_n-x_i).
\]

Cada factor del producto es distinto de cero porque las abscisas son distintas. Por eso se puede despejar de forma única:

\[
\boxed{a_n=
\frac{y_n-p_{n-1}(x_n)}{\prod_{i=0}^{n-1}(x_n-x_i)}},
\]

\[
\boxed{p_n(x)=p_{n-1}(x)+a_n\prod_{i=0}^{n-1}(x-x_i)}.
\]

La fórmula construye un candidato que pasa por todos los puntos, lo que da existencia. Al mismo tiempo, todo candidato estaba obligado a tener ese $q$ y ese único $a_n$, lo que da unicidad. No se necesita una segunda prueba independiente con un sistema matricial.

Geométricamente, de una constante se pasa a una recta agregando una función lineal que vale cero en el primer nodo. De una recta se pasa a una parábola agregando un múltiplo de $(x-x_0)(x-x_1)$: se mantienen los dos primeros puntos y se ajusta la altura del tercero. Esa familia de correcciones es la que el docente dibuja para explicar la inducción.

\begin{figure}[H]\centering
\begin{tikzpicture}
\begin{axis}[width=10.5cm,height=6cm,axis lines=left,xmin=-0.1,xmax=2.2,ymin=0,ymax=3.3,xtick={0,1,2},xticklabels={$x_0$,$x_1$,$x_2$},ytick=\empty,xlabel={$x$},clip=false]
\addplot[figgray,dashed,domain=0:2.1,forget plot]{1};
\addplot[figblue,thick,domain=0:2.1,forget plot]{1+x*(x-1)};
\addplot[figamber,densely dashdotted,line width=1pt,domain=0:2.1,forget plot]{1+0.5*x*(x-1)};
\addplot[only marks,mark=*,mark options={fill=figred,draw=figred},forget plot] coordinates {(0,1)(1,1)(2,3)};
\node[anchor=east,font=\small] at (axis cs:1.95,3.05){$(x_2,y_2)$};
\node[anchor=west,font=\small,figgray] at (axis cs:2.12,1){$p_1$};
\node[anchor=west,font=\small,figblue] at (axis cs:2.12,3.31){$p_2$};
\node[anchor=west,font=\small,figamber] at (axis cs:2.12,2.155){otro $a_2$};
\end{axis}
\end{tikzpicture}
\caption{Mecanismo de la inducción: sumar $a_2(x-x_0)(x-x_1)$ mantiene los dos primeros valores y permite ajustar el tercero. Coordenadas ilustrativas elegidas para representar la familia descrita, no datos numéricos del pizarrón.}
\end{figure}

El argumento es constructivo, pero en esta clase no se desarrolla todavía su implementación como fórmula de Newton. Lo que sigue es otra representación del mismo polinomio, mediante la base monomial.

\section{El sistema de Vandermonde}

\subsection{De condiciones puntuales a ecuaciones en los coeficientes}

Se escribe

\[
p_n(x)=c_0+c_1x+c_2x^2+\cdots+c_nx^n.
\]

Al evaluar en cada nodo, aparecen $n+1$ ecuaciones:

\[
\begin{aligned}
c_0+c_1x_0+c_2x_0^2+\cdots+c_nx_0^n&=y_0,\\
c_0+c_1x_1+c_2x_1^2+\cdots+c_nx_1^n&=y_1,\\
&\vdots\\
c_0+c_1x_n+c_2x_n^2+\cdots+c_nx_n^n&=y_n.
\end{aligned}
\]

Los nodos y las alturas son datos. Las incógnitas son los coeficientes $c_j$. Aunque se usan potencias de las abscisas, el sistema es lineal en esas incógnitas:

\[
\underbrace{\begin{pmatrix}
1&x_0&x_0^2&\cdots&x_0^n\\
1&x_1&x_1^2&\cdots&x_1^n\\
\vdots&\vdots&\vdots&&\vdots\\
1&x_n&x_n^2&\cdots&x_n^n
\end{pmatrix}}_{V}
\underbrace{\begin{pmatrix}c_0\\c_1\\\vdots\\c_n\end{pmatrix}}_{\mathbf c}
=\underbrace{\begin{pmatrix}y_0\\y_1\\\vdots\\y_n\end{pmatrix}}_{\mathbf y}.
\]

Esta es la \textbf{matriz de Vandermonde}, de tamaño $(n+1)\times(n+1)$. Cada fila evalúa todos los monomios en un nodo; cada columna corresponde a una potencia fija. Su invertibilidad se deduce de la existencia y unicidad ya probadas, sin calcular su determinante. Si se repitieran abscisas, se repetirían filas en esta convención.

\subsection{Orden de columnas en la computadora}

El ejemplo usa tres abscisas $1/4,1,5/2$. En el orden creciente de potencias, la matriz es

\[
V=\begin{pmatrix}
1&1/4&1/16\\
1&1&1\\
1&5/2&25/4
\end{pmatrix}.
\]

La función \texttt{vander} mostrada ordena las columnas al revés, empezando por la mayor potencia:

\[
V_{\mathrm{desc}}=\begin{pmatrix}
1/16&1/4&1\\
1&1&1\\
25/4&5/2&1
\end{pmatrix},\qquad
V_{\mathrm{desc}}\begin{pmatrix}c_2\\c_1\\c_0\end{pmatrix}
=\begin{pmatrix}y_0\\y_1\\y_2\end{pmatrix}.
\]

Resolver con la contrabarra entrega los coeficientes en ese orden. Para escribir o evaluar el polinomio hay que conservar la correspondencia entre posición y potencia; invertir sólo las columnas, sin reinterpretar los coeficientes, sería otro polinomio.

\begin{algorithm}[H]
\small\DontPrintSemicolon
\hspace*{0.0em}\texttt{Entrada: nodos distintos x[0:n], alturas y[0:n]}\;
\hspace*{0.0em}\texttt{Construir V[i,j] = x[i]\textasciicircum{}j, con V[i,0] = 1}\;
\hspace*{0.0em}\texttt{Resolver V*c = y}\;
\hspace*{0.0em}\texttt{Evaluar p(t) = c[0] + c[1]*t + ... + c[n]*t\textasciicircum{}n}\;
\hspace*{0.0em}\texttt{Salida: coeficientes y evaluaciones del interpolante}\;
\end{algorithm}

El esquema usa potencias crecientes, a diferencia de \texttt{vander}. Cada fila asegura una condición de interpolación; resolver obtiene todos los coeficientes conjuntamente y evaluar es una tarea posterior. Se añade explícitamente $V[i,0]=1$ para no confundir la columna constante cuando algún nodo es cero.

\begin{notebox}[Alcance y precisiones]
El audio no enumera las ordenadas de este ejemplo y vacila al leer los coeficientes. Se recuperan las abscisas, las dos matrices y el procedimiento, pero no se inventan la parábola ni sus coeficientes exactos.
\end{notebox}

\section{Condicionamiento y cambios en los datos}

\subsection{Una matriz invertible puede ser numéricamente difícil}

La facilidad conceptual de Vandermonde no garantiza buena precisión. La demostración arma nodos equiespaciados en $[0,1]$, aumenta su cantidad y estima condiciones. El gráfico logarítmico muestra un crecimiento fuerte; para veinte puntos se anuncia aproximadamente $4\cdot10^{16}$.

Ese factor puede destruir el control de error en los coeficientes aunque el sistema se resuelva con un procedimiento que produzca residuo pequeño. En $V\mathbf c=\mathbf y$, el residuo describe el desacuerdo al evaluar en los nodos. No controla por sí solo las evaluaciones entre ellos ni la precisión de cada coeficiente cuando $V$ está mal condicionada.

\begin{notebox}[Alcance y precisiones]
La transcripción no contiene la tabla completa del gráfico de condición. Se registra el procedimiento y el valor anunciado, sin dibujar una curva numérica ficticia que parezca una medición de la clase.
\end{notebox}

\subsection{Interpolar la constante uno}

Se toman nodos en $[0,1]$ y todas las alturas iguales a uno. Por unicidad, el interpolante exacto es $p(x)=1$, independientemente de cuántos nodos distintos se usen. En el orden creciente de potencias, esto significa

\[
V\begin{pmatrix}1\\0\\\vdots\\0\end{pmatrix}
=\begin{pmatrix}1\\1\\\vdots\\1\end{pmatrix}.
\]

Con veinticinco puntos, la curva calculada parece próxima a uno, pero sus coeficientes ya son extraños. El docente aumenta a cincuenta y cinco puntos y observa oscilaciones grandes, valores del gráfico cercanos a nueve y coeficientes del orden de $10^{15}$. La computadora advierte que el sistema es muy problemático. No se conserva el perfil exacto de esas oscilaciones porque el audio no lo determina.

El ejemplo separa el problema matemático, cuya respuesta es una constante trivial, de la representación elegida para calcularla. Resolver un sistema grande de Vandermonde puede ser una mala manera de encontrar un objeto sencillo. Tener residuos razonables en los nodos no equivale a haber obtenido coeficientes fiables.

\subsection{Agregar o corregir datos}

Si se incorpora otro punto, cambian el tamaño de la matriz, el vector de datos y la cota de grado. El procedimiento monomial presentado exige armar y resolver de nuevo el sistema ampliado. Si cambia solamente una ordenada, la matriz permanece igual, pero hay que resolver con el nuevo lado derecho; no se cambia directamente un único coeficiente del polinomio.

La clase siguiente buscará otras dos formas de representar el mismo interpolante para evitar el sistema mal condicionado y facilitar cambios en los datos. No se abandona la unicidad del polinomio: se cambia la forma de construirlo y expresarlo.
\end{document}
